\subsection{Image of a complex measure}

Let \(\theta\) be a complex measure on T, and let \(\pi\) be a mapping of T into a locally compact space X; suppose that \(\pi\) is \(\theta\) -measurable, and that for each \(f \in \mathcal{K}(X; \mathbf{C})\) , \(f \circ \pi\) is essentially \(\theta\) -integrable. Since it is equivalent to say that a function is measurable (resp. essentially integrable) with respect to \(\theta\) or with respect to \(|\theta|\) , this means that \(\pi\) is \(|\theta|\) -proper. If \(f \in \mathcal{K}(X; \mathbf{C})\) ,

\[
\left| \int (f \circ \pi) d \theta \right| \leqslant \int (| f | \circ \pi) d | \theta |;\tag{7}
\]

it follows at once that the linear form \(f \mapsto \int(f \circ \pi) d\theta\) on \(\mathcal{K}(X; \mathbf{C})\) is a complex measure on X (Ch. III, §1, No.~3, Prop. 6), and one can make the following definition:

DEFINITION 2. --- Let \(\theta\) be a complex measure on a locally compact space T. A mapping \(\pi\) of T into a locally compact space X is said to be \(\theta\) -proper if it is \(|\theta|\) -proper. The measure \(f \mapsto \int(f \circ \pi) d\theta\) is then called the image of \(\theta\) under \(\pi\) and is denoted \(\pi(\theta)\) .

The relation (7) may then be put in the following form:

\[
| \pi (\theta) | \leqslant \pi (| \theta |).\tag{8}
\]

The measure \(\pi(\theta)\) may be zero without \(\theta\) being zero, as one sees immediately on taking for T a space reduced to two points a, b, for \(\theta\) the measure \(\varepsilon_{a} - \varepsilon_{b}\) , and for \(\pi\) a constant mapping.

Let \(\theta\) and \(\theta'\) be two complex measures on T; if \(\pi\) is \(\theta\) -proper and \(\theta'\) -proper, it follows from Cor. 2 of Prop. 6 that \(\pi\) is \((\theta + \theta')\) -proper since \(|\theta + \theta'| \leqslant |\theta| + |\theta'|\) , and obviously \(\pi(\theta + \theta') = \pi(\theta) + \pi(\theta')\) .

In particular, if \(\theta\) is a real measure and \(\pi\) is \(\theta\) -proper, then

\[
\pi (\theta) = \pi (\theta^ {+}) - \pi (\theta^ {-}).\tag{9}
\]

Several results established earlier extend at once to complex measures; we cite the most important among them.

PROPOSITION 7. --- Let \(\theta\) be a complex measure on T, \(\pi\) a \(\theta\) -proper mapping of T into a locally compact space X, \(\nu\) the image measure \(\pi(\theta)\) .

\begin{enumerate}
\def\labelenumi{\alph{enumi})}
\item
  Let A be a subset of X; if \(\overline{\pi}^{1}(A)\) is locally \(\theta\) -negligible, then A is locally \(\nu\) -negligible.
\item
  Let \(f\) be a mapping of \(X\) into a topological space; if \(f \circ \pi\) is \(\theta\) -measurable, then \(f\) is \(\nu\) -measurable.
\item
  Let \(\mathbf{f}\) be a function defined on \(X\) , with values in a Banach space \(F\) ; if \(\mathbf{f} \circ \pi\) is essentially \(\theta\) -integrable, then \(\mathbf{f}\) is essentially \(\nu\) -integrable and
\end{enumerate}

\[
\int \mathbf {f} (\pi (t)) d \theta (t) = \int \mathbf {f} (x) d \nu (x).\tag{10}
\]

Taking into account formula (8), these results may be deduced from Cor. 2 of Prop. 2, from Prop. 3, and from Th. 1 of No.~2.

